% Figure 3 -- the ten implemented evidence bases, grouped by the strength they can carry.
% Names, conditions and strengths are taken verbatim from EVIDENCE_ARBITRATION_SPEC.md section 3
% (and evarb/arbiter.py). No basis is renamed, added or omitted.
\begin{figure}[t]
\centering
\begin{tikzpicture}[font=\footnotesize, x=1mm, y=1mm,
  base/.style={draw=cBoxEdge, line width=0.4pt, rounded corners=1.2pt, align=left,
               inner sep=2.4pt, text width=46mm, font=\scriptsize, text=cInk, minimum height=9mm},
  hdr/.style={align=center, font=\scriptsize\bfseries, text=white, inner sep=2.2pt,
              rounded corners=1.2pt, minimum height=5.2mm, text width=48mm}]

% ---- column headers: the gate admits the left column only
\node[hdr, fill=cEBG]   (hs) at (26,86) {STRONG\quad\textnormal{\scriptsize the only column the gate admits}};
\node[hdr, fill=cAll, text=cInk] (hm) at (80,86) {MODERATE\quad\textnormal{\scriptsize enforced only by \cond{EBG-all}}};
\node[hdr, fill=cNone]  (hn) at (134,86) {NONE\quad\textnormal{\scriptsize never enforced}};

% ---- strong
\node[base, fill=cAccentFill] (b1) at (26,78)
  {\textbf{\cond{explicit\_correction}}\\asserted correction or dismissal referring to the rival, at different dates};
\node[base, fill=cAccentFill] (b2) at (26,65.5)
  {\textbf{\cond{explicit\_supersession}}\\asserted replacement or revision referring to the rival, at different dates};
\node[base, fill=cAccentFill] (b3) at (26,53)
  {\textbf{\cond{same\_day\_correction}}\\asserted correction, dismissal or supersession between same-date records, referring to the rival};
\node[base, fill=cAccentFill] (b4) at (26,40.5)
  {\textbf{\cond{valid\_time}}\\every established record states a validity date; the latest is unique};
\node[base, fill=cAccentFill] (b5) at (26,28)
  {\textbf{\cond{record\_time}}\\unique latest \cond{recorded\_on}};
\node[base, fill=cAccentFill] (b6) at (26,15.5)
  {\textbf{\cond{genuine\_contradiction}}\\same time, different values, no ordering evidence, no cue of any kind};
\node[base, fill=cAccentFill] (b7) at (26,3)
  {\textbf{\cond{negated\_correction}}\\negated correction with an asserted ``both stand / disagree'' marker};

% ---- moderate (demotions of the same bases, plus two of their own)
\node[base, fill=cWarnFill] (m1) at (80,78)
  {\textbf{\cond{explicit\_correction} / \cond{explicit\_supersession}}\\\emph{demoted} when the correcting record is the older one};
\node[base, fill=cWarnFill] (m2) at (80,65.5)
  {\textbf{\cond{same\_day\_correction}}\\\emph{demoted} when the correction is unreferenced and only breaks a same-time tie};
\node[base, fill=cWarnFill] (m3) at (80,53)
  {\textbf{\cond{record\_time}}\\\emph{demoted} when the winner's own text signals doubt (``proposed'', ``planned'', \ldots)};
\node[base, fill=cWarnFill] (m4) at (80,40.5)
  {\textbf{\cond{genuine\_contradiction}}\\\emph{demoted} when any cue is present};
\node[base, fill=cWarnFill] (m5) at (80,28)
  {\textbf{\cond{negated\_correction}}\\negated replacement without a marker, retention of the rival, or an unconfirmed change: that record is set aside};
\node[base, fill=cWarnFill] (m6) at (80,15.5)
  {\textbf{\cond{source\_authority}}\\same-time tie in which exactly one value comes from an authoritative source};
\node[base, fill=cWarnFill] (m7) at (80,3)
  {\textbf{\cond{quoted\_or\_hypothetical}}\\the record's own change or value is hypothetical: that record is set aside};

% ---- none
\node[base, fill=white] (n1) at (134,78)
  {\textbf{\cond{quoted\_or\_hypothetical}}\\the only ordering cue is reported speech $\rightarrow$ \cond{unresolved}};
\node[base, fill=white] (n2) at (134,65.5)
  {\textbf{\cond{insufficient\_evidence}}\\missing dates without text dates; cyclic or inconsistent cues $\rightarrow$ \cond{unresolved}};

\node[base, draw=none, fill=none, text width=50mm] (note) at (134,48)
  {\textbf{Ten bases, three strengths.} Six bases can carry more than one strength depending on their
   conditions, which is why gating cannot be replaced by a fixed list of ``trusted'' cue types: the same
   cue licenses enforcement in one context and only advice in another.};

\node[base, draw=cEnf, fill=cWarnFill, text width=50mm] (note2) at (134,22)
  {\textbf{What the gate costs.} Moving the gate from \emph{strong} to \emph{strong\,+\,moderate} ---
   that is, enabling the middle column --- raised task success by $9.7$\pp{} and the harmful-answer rate from
   3.1\% to 6.9\% (\Cref{sec:ablations}).};

\draw[decorate, decoration={brace, amplitude=3.5pt}, draw=cEBG, line width=0.6pt]
  (51.8,83.5) -- (51.8,-2.5);
\end{tikzpicture}

\vspace{1mm}
\caption[The ten evidence bases of Evarb-0.2]{%
\textbf{The ten evidence bases, grouped by the strength they can carry.} Names and conditions are verbatim from
the frozen specification; no basis is renamed, added or omitted. The left column is what gated enforcement acts on.
Six of the ten bases appear in more than one column because their strength depends on their conditions --- a
correction whose correcting record is the older one is demoted, an ordering by record date is demoted when the
winning record's own text signals doubt, and a contradiction is demoted when any cue is present. This
context-dependence is the reason the mechanism reports a basis and a strength rather than a verdict.\\
\emph{Accessible description:} a three-column layout headed STRONG, MODERATE and NONE, holding seven, seven and two
entries respectively; a brace marks the left column as the only one the gate admits.\\
\emph{Source:} \srcfile{EVIDENCE_ARBITRATION_SPEC.md} \S3; \srcfile{evarb/arbiter.py}.}
\label{fig:bases}
\end{figure}
